\section{Position Hold PID}

\begin{figure}[h]
\centering
\resizebox{\textwidth}{!}{
\begin{tikzpicture}[
    block/.style={
        rectangle,
        draw,
        rounded corners,
        align=center,
        minimum width=5cm,
        minimum height=1.2cm
    },
    arrow/.style={
        -{Latex},
        thick
    },
    node distance=4cm
]

\node (input)
{};

\node[block, right=of input] (pos)
{Position PID};

\node[right=of pos] (output)
{};

\draw[arrow] (input) --
node[midway, above] {$x_d,y_d$}
node[midway, below] {$x,y,v_x,v_y,\psi$}
(pos);

\draw[arrow] (pos) -- node[midway, above]
{$\phi_d,\theta_d$} (output);

\end{tikzpicture}
}
\caption{Position PID control law}
\label{fig:position_pid}
\end{figure}


\subsection{Công thức PID chuẩn}

Công thức PID tổng quát có dạng \cite{astrom1995pid}:

\begin{equation}
    u(t) = K_p e(t) + K_i \int e(t)\,dt + K_d \dot{e}(t)
\end{equation}

trong đó:

\begin{equation}
    e(t) = r(t) - y(t)
\end{equation}

Với:
\begin{itemize}
    \item $r(t)$: giá trị mong muốn,
    \item $y(t)$: giá trị đo được,
    \item $e(t)$: sai số,
    \item $u(t)$: output điều khiển,
    \item $K_p, K_i, K_d$: hệ số tỉ lệ, tích phân và vi phân.
\end{itemize}

Trong bộ điều khiển drone, output $u(t)$ không nhất thiết là PWM.
Tùy tầng điều khiển, $u(t)$ có thể là gia tốc, góc mong muốn, moment,
lực nâng hoặc PWM.

\subsection{Sai số vị trí trong Position Hold}

Gọi vị trí mong muốn của drone trong WORLD-coordination [W] là:

\begin{equation}
    x_d,\quad y_d
\end{equation}

Vị trí hiện tại của drone là:

\begin{equation}
    x,\quad y
\end{equation}

Sai số vị trí được định nghĩa là:

\begin{equation}
    e_x = x_d - x
\end{equation}

\begin{equation}
    e_y = y_d - y
\end{equation}

Nếu $e_x > 0$, drone đang nằm phía sau target theo trục $x$ và cần tạo gia tốc
theo chiều $+x$.

Nếu $e_y > 0$, drone đang lệch so với target theo trục $y$ và cần tạo gia tốc
theo chiều $+y$.

\subsection{Sai số vận tốc}

Gọi vận tốc mong muốn là:

\begin{equation}
    v_{xd},\quad v_{yd}
\end{equation}

Vận tốc hiện tại là:

\begin{equation}
    v_x,\quad v_y
\end{equation}

Sai số vận tốc là:

\begin{equation}
    e_{v_x} = v_{xd} - v_x
\end{equation}

\begin{equation}
    e_{v_y} = v_{yd} - v_y
\end{equation}

Ta có:

\begin{equation}
    e_x = x_d - x
\end{equation}

Lấy đạo hàm hai vế theo thời gian:

\begin{equation}
    \dot{e}_x = \dot{x}_d - \dot{x}
\end{equation}

Mà:

\begin{equation}
    \dot{x}_d = v_{xd}
\end{equation}

\begin{equation}
    \dot{x} = v_x
\end{equation}

Do đó:

\begin{equation}
    \dot{e}_x = v_{xd} - v_x = e_{v_x}
\end{equation}

Tương tự:

\begin{equation}
    \dot{e}_y = v_{yd} - v_y = e_{v_y}
\end{equation}

Vì vậy, thành phần sai số vận tốc $e_{v_x}$ và $e_{v_y}$ chính là thành phần
vi phân của sai số vị trí.

\subsection{Tích phân sai số}

Thành phần tích phân của sai số vị trí là:

\begin{equation}
    I_x = \int e_x\,dt
\end{equation}

\begin{equation}
    I_y = \int e_y\,dt
\end{equation}

Trong dạng rời rạc:

\begin{equation}
    I_x(k) = I_x(k-1) + e_x(k)\Delta t
\end{equation}

\begin{equation}
    I_y(k) = I_y(k-1) + e_y(k)\Delta t
\end{equation}

Để tránh hiện tượng tích phân tăng quá lớn, ta giới hạn:

\begin{equation}
    -I_{\max} \leq I_x \leq I_{\max}
\end{equation}

\begin{equation}
    -I_{\max} \leq I_y \leq I_{\max}
\end{equation}

Giới hạn này giúp chống hiện tượng \textit{integral windup}.

\subsection{PID vị trí tạo gia tốc mong muốn}

Trong tầng Position Hold, PID không tạo PWM trực tiếp.
Output của PID vị trí là gia tốc ngang mong muốn:

\begin{equation}
    a_{x,cmd}
\end{equation}

\begin{equation}
    a_{y,cmd}
\end{equation}

Công thức PID theo trục $x$:

\begin{equation}
    a_{x,cmd}
    =
    K_{p_x} e_x
    +
    K_{d_x} e_{v_x}
    +
    K_{i_x} I_x
\end{equation}

Công thức PID theo trục $y$:

\begin{equation}
    a_{y,cmd}
    =
    K_{p_y} e_y
    +
    K_{d_y} e_{v_y}
    +
    K_{i_y} I_y
\end{equation}

So với công thức PID chuẩn:

\begin{equation}
    u = K_p e + K_i \int e\,dt + K_d \dot{e}
\end{equation}

thì trong tầng Position Hold:

\begin{equation}
    u_x = a_{x,cmd}
\end{equation}

\begin{equation}
    u_y = a_{y,cmd}
\end{equation}

Tức là output của PID vị trí là gia tốc ngang mong muốn, không phải PWM.

\subsection{Bản chất vật lý của việc chọn gia tốc làm output}

Theo định luật II Newton:

\begin{equation}
    F = ma
\end{equation}

Muốn thay đổi vị trí ngang của drone, ta cần tạo gia tốc ngang:

\begin{equation}
    a_x,\quad a_y
\end{equation}

Do đó, tầng Position Hold trả lời câu hỏi:

\begin{center}
    Drone cần gia tốc ngang bao nhiêu để quay về vị trí mong muốn?
\end{center}

Vì vậy, output hợp lý của tầng Position Hold là:

\begin{equation}
    a_{x,cmd},\quad a_{y,cmd}
\end{equation}

Sau đó, các gia tốc này sẽ được chuyển thành góc nghiêng mong muốn.

\subsection{Giới hạn gia tốc ngang}

Vector gia tốc ngang mong muốn là:

\begin{equation}
    \mathbf{a}_{cmd}
    =
    \begin{bmatrix}
        a_{x,cmd} \\
        a_{y,cmd}
    \end{bmatrix}
\end{equation}

Độ lớn vector gia tốc là:

\begin{equation}
    \|\mathbf{a}_{cmd}\|
    =
    \sqrt{a_{x,cmd}^{2}+a_{y,cmd}^{2}}
\end{equation}

Project giới hạn gia tốc ngang theo điều kiện:

\begin{equation}
    \|\mathbf{a}_{cmd}\| \leq a_{\max}
\end{equation}

Nếu:

\begin{equation}
    \sqrt{a_{x,cmd}^{2}+a_{y,cmd}^{2}} > a_{\max}
\end{equation}

thì scale lại:

\begin{equation}
    a_{x,cmd}^{new}
    =
    a_{x,cmd}
    \frac{a_{\max}}
    {\sqrt{a_{x,cmd}^{2}+a_{y,cmd}^{2}}}
\end{equation}

\begin{equation}
    a_{y,cmd}^{new}
    =
    a_{y,cmd}
    \frac{a_{\max}}
    {\sqrt{a_{x,cmd}^{2}+a_{y,cmd}^{2}}}
\end{equation}

Bản chất của bước này là giới hạn độ mạnh của lệnh điều khiển vị trí.
Nếu gia tốc ngang quá lớn, drone phải nghiêng quá nhiều, dễ gây mất ổn định.

\subsection{Đổi từ [W] sang [B]}

Gia tốc $a_{x,cmd}$ và $a_{y,cmd}$ đang nằm trong [W].
Tuy nhiên, roll và pitch lại là góc nghiêng theo thân drone.

Gọi yaw hiện tại của drone là:

\begin{equation}
    \psi
\end{equation}

Vector gia tốc trong [W] là:

\begin{equation}
    \mathbf{a}^{W}
    =
    \begin{bmatrix}
        a_x \\
        a_y
    \end{bmatrix}
\end{equation}

Trục forward của drone trong [W] là:

\begin{equation}
    \mathbf{e}_f
    =
    \begin{bmatrix}
        \cos\psi \\
        \sin\psi
    \end{bmatrix}
\end{equation}

Trục left của drone trong [W] là:

\begin{equation}
    \mathbf{e}_l
    =
    \begin{bmatrix}
        -\sin\psi \\
        \cos\psi
    \end{bmatrix}
\end{equation}

Gia tốc theo hướng forward của drone là hình chiếu của vector gia tốc world
lên trục forward:

\begin{equation}
    a_f
    =
    \mathbf{a}^{W} \cdot \mathbf{e}_f
\end{equation}

Suy ra:

\begin{equation}
    a_f
    =
    a_x\cos\psi + a_y\sin\psi
\end{equation}

Gia tốc theo hướng left của drone là:

\begin{equation}
    a_l
    =
    \mathbf{a}^{W} \cdot \mathbf{e}_l
\end{equation}

Suy ra:

\begin{equation}
    a_l
    =
    -a_x\sin\psi + a_y\cos\psi
\end{equation}

Trong đó:
\begin{itemize}
    \item $a_f$: gia tốc mong muốn theo hướng tiến/lùi của drone,
    \item $a_l$: gia tốc mong muốn theo hướng trái/phải của drone.
\end{itemize}

\subsection{Bản chất vật lý - nghiêng drone để tạo gia tốc ngang}

Quadcopter tạo lực nâng tổng $T$ theo trục thân của nó \cite{bouabdallah2007design}.

Khi drone nằm ngang, lực nâng hướng thẳng lên.
Ở trạng thái hover:

\begin{equation}
    T \approx mg
\end{equation}

Khi drone nghiêng một góc $\alpha$, lực nâng bị nghiêng theo.
Khi đó lực nâng có thành phần đứng và thành phần ngang:

\begin{equation}
    T_z = T\cos\alpha
\end{equation}

\begin{equation}
    T_h = T\sin\alpha
\end{equation}

Để giữ độ cao gần như không đổi:

\begin{equation}
    T\cos\alpha \approx mg
\end{equation}

Theo phương ngang:

\begin{equation}
    T\sin\alpha = ma_h
\end{equation}

Chia hai phương trình:

\begin{equation}
    \frac{T\sin\alpha}{T\cos\alpha}
    =
    \frac{ma_h}{mg}
\end{equation}

Suy ra:

\begin{equation}
    \tan\alpha
    =
    \frac{a_h}{g}
\end{equation}

Do đó:

\begin{equation}
    \alpha
    =
    \arctan\left(\frac{a_h}{g}\right)
\end{equation}

Đây là công thức vật lý cơ bản để chuyển từ gia tốc ngang mong muốn sang
góc nghiêng mong muốn.

\subsection{Tính pitch target}

Gọi gia tốc mong muốn theo hướng forward là:

\begin{equation}
    a_f
\end{equation}

Góc pitch mong muốn là:

\begin{equation}
    \theta_d
\end{equation}

Theo quan hệ vật lý:

\begin{equation}
    \tan\theta_d
    =
    \frac{a_f}{g}
\end{equation}

Suy ra:

\begin{equation}
    \theta_d
    =
    \arctan\left(\frac{a_f}{g}\right)
\end{equation}

Trong project, công thức được viết tương đương dưới dạng:

\begin{equation}
    \theta_d
    =
    \operatorname{atan2}(a_f, g)
\end{equation}

Với góc nhỏ:

\begin{equation}
    \theta_d \approx \frac{a_f}{g}
\end{equation}

Trong đó:

\begin{equation}
    \theta_d = pitch\_target
\end{equation}

\subsection{Tính roll target}

Gọi gia tốc mong muốn theo hướng left là:

\begin{equation}
    a_l
\end{equation}

Góc roll mong muốn là:

\begin{equation}
    \phi_d
\end{equation}

Theo quy ước dấu trong project:

\begin{equation}
    \phi_d
    =
    -\operatorname{atan2}(a_l, g)
\end{equation}

Tương đương:

\begin{equation}
    \phi_d
    =
    -\arctan\left(\frac{a_l}{g}\right)
\end{equation}

Với góc nhỏ:

\begin{equation}
    \phi_d \approx -\frac{a_l}{g}
\end{equation}

Trong đó:

\begin{equation}
    \phi_d = roll\_target
\end{equation}

Dấu âm xuất hiện do quy ước trục roll trong mô hình.
Nếu sau này drone bay ngược hướng trái/phải, đây là một dấu cần kiểm tra.

\subsection{Giới hạn góc nghiêng}

Sau khi tính được $\phi_d$ và $\theta_d$, project giới hạn:

\begin{equation}
    -\alpha_{\max}
    \leq
    \phi_d
    \leq
    \alpha_{\max}
\end{equation}

\begin{equation}
    -\alpha_{\max}
    \leq
    \theta_d
    \leq
    \alpha_{\max}
\end{equation}

Nếu:

\begin{equation}
    \alpha_{\max} = 6.5^\circ
\end{equation}

thì gia tốc ngang tối đa tương ứng xấp xỉ là:

\begin{equation}
    a_{\max}
    =
    g\tan(\alpha_{\max})
\end{equation}

Thay số:

\begin{equation}
    a_{\max}
    =
    9.81\tan(6.5^\circ)
\end{equation}

\begin{equation}
    a_{\max}
    \approx
    1.12\ \mathrm{m/s^2}
\end{equation}

Điều này cho thấy giới hạn góc nghiêng và giới hạn gia tốc ngang có liên hệ
trực tiếp với nhau.

\subsection{Công thức tổng hợp của Position Hold PID}

Toàn bộ tầng Position Hold có thể viết tóm tắt như sau.

Sai số vị trí:

\begin{equation}
    e_x = x_d - x
\end{equation}

\begin{equation}
    e_y = y_d - y
\end{equation}

Sai số vận tốc:

\begin{equation}
    e_{v_x} = v_{xd} - v_x
\end{equation}

\begin{equation}
    e_{v_y} = v_{yd} - v_y
\end{equation}

Tích phân sai số:

\begin{equation}
    I_x = \int e_x\,dt
\end{equation}

\begin{equation}
    I_y = \int e_y\,dt
\end{equation}

PID tạo gia tốc ngang mong muốn:

\begin{equation}
    a_x
    =
    K_{p_x}e_x
    +
    K_{d_x}e_{v_x}
    +
    K_{i_x}I_x
\end{equation}

\begin{equation}
    a_y
    =
    K_{p_y}e_y
    +
    K_{d_y}e_{v_y}
    +
    K_{i_y}I_y
\end{equation}

Đổi sang hệ trục của drone:

\begin{equation}
    a_f
    =
    a_x\cos\psi + a_y\sin\psi
\end{equation}

\begin{equation}
    a_l
    =
    -a_x\sin\psi + a_y\cos\psi
\end{equation}

Đổi gia tốc thành góc nghiêng:

\begin{equation}
    \theta_d
    =
    \operatorname{atan2}(a_f, g)
\end{equation}

\begin{equation}
    \phi_d
    =
    -\operatorname{atan2}(a_l, g)
\end{equation}

\vspace{0.5cm}
\hrule

